\section{Scope and Validity Considerations}
\label{sec:limitations}

The study uses \dataset{} as a realistic synthetic benchmark with explicit burst, congestion, outage, degradation, and maintenance windows. Race-condition pressure is represented through precomputed contention and queue carryover instead of live mutable queues. This choice keeps the comparison repeatable and the evaluation auditable.

The compared methods are implemented inside the same evaluation pipeline for comparison fairness. They share the same features, temporal split, and metrics, which keeps the evaluation aligned across policies. The FL-DDPG-style method serves as the federated actor-critic comparator within this common pipeline.

The multi-seed analysis uses three seeds. It provides a repeatability check for the proposed configuration and ablations. Future extensions can add more seeds, confidence intervals for all baselines, independent workload-generator samples, and a no-bootstrapping ablation with $\gamma=0$ to isolate the value of temporal bootstrapping. No unrun baseline-variance or $\gamma=0$ results are included in the present tables.

Federated training is privacy-aligned because raw zone data are not pooled in the algorithmic design. Secure aggregation, differential privacy, and attack evaluation against model updates remain outside the present implementation and are natural extensions for a security-focused study. Communication efficiency is therefore treated as configuration-dependent; the present manuscript reports the transfer expression but does not claim that model-update traffic is always smaller than centralizing all raw features.

Oracle and supervised label generation depend on the same latency functions and state tables used by evaluation. This keeps the benchmark internally consistent and makes Oracle a useful generated latency reference. Measured-trace studies can extend the same evaluation protocol with deployment-specific overheads and online queue mutation.

The pipeline evaluates offloading decisions in an offline replay style. The dataset contains precomputed contention, while a future live simulator could allow each policy action to mutate subsequent queues, bandwidth pressure, and admission state. This would extend the benchmark into an online closed-loop evaluation setting and would test whether the same reward and validation design remains stable when actions change future states rather than only current-task rewards.
